\begin{sloppypar}
, Slotted Iris Wakefields, and FEL InteractionThe linear accelerating infrastructure utilizes $5.712\text{ GHz}$ constant-gradient C-band accelerating structures operating in the $2\pi/3$ traveling-wave mode at an accelerating gradient $E_{\text{acc}} = 38.5\text{ MV/m}$. Continuous 1D steady-state approximations are replaced by an intra-burst train of $N_b = 2,500$ discrete micro-bunches injected at the fundamental RF period $\tau_b = 175.07\text{ ps}$ with an intra-burst macroscopic current $I_{\text{burst}} = 1.1424\text{ A}$. The discrete charge per micro-bunch is $q_b = 0.200\text{ nC}$ ($1.2483 \times 10^9$ electrons/bunch), producing a total macro-pulse duration of $437.675\text{ ns}$.Transverse beam stability requires active attenuation of dipole higher-order mode passbands centered at $8.980\text{ GHz}$ and $11.424\text{ GHz}$. Each cell iris contains four symmetric radial sectoral slots that couple dipole wall currents into radial waveguides terminated with silicon-carbide dielectric absorbers. This geometry achieves an external quality factor $Q_{\text{ext}} \le 50.0$ for all deflecting modes while preserving the fundamental accelerating mode quality factor ($Q_0 = 11,500$) and shunt impedance ($R_\parallel = 82\text{ M}\Omega/\text{m}$) within $98.8\%$. The wakefield damping time is $\tau_d = 1.772\text{ ns}$ ($10.12$ micro-bunch periods), ensuring wakefield memory attenuation below $1.0\%$ within 46 bunches.Cumulative multi-bunch beam breakup across the 120-meter linac is governed by the time-domain discrete bunch-by-bunch tracking formulation:\[
\resizebox{\columnwidth}{!}{$\displaystyle \frac{d}{ds} \left[ \gamma_M(s) \frac{d x_M(s)}{ds} \right] + \gamma_M(s) k_\beta^2(s) x_M(s) = \frac{e^2 N_e}{m_e c^2} \sum_{k=1}^{M-1} x_k(s) W_\perp\left((M - k)\tau_b\right)$}
\]
Incorporating cell manufacturing tolerances ($\pm 2\,\mu\text{m}$) and an uncorrelated dipole detuning spread of $\sigma_\omega / \bar{\omega} = 0.15\%$ introduces phase cancellation into the wakefield sum. Numerical integration demonstrates that the cumulative BBU amplification factor is limited to $A_{\text{BBU}} \le 1.184$, guaranteeing that the projected transverse emittance satisfies $\epsilon_{n,x} = 0.442\text{ mm}\cdot\text{mrad} \le 0.50\text{ mm}\cdot\text{mrad}$ and the intra-burst energy spread remains bounded at $\Delta \gamma / \gamma = 8.65 \times 10^{-5} \le 1.00 \times 10^{-4}$ at the target interface.Downstream micro-bunch conditioning employs an optical klystron FEL composed of a modulator undulator, an achromatic dispersion chicane ($R_{56}$), and a radiator tuned to the 5th harmonic ($h = 5$) of a $266\text{ nm}$ seed laser. Operating with a bunching parameter $D = 1.34$, the system achieves an optical harmonic micro-bunching factor $b_5 = 0.284$. Focused to a diffraction-limited spot radius $w_0 = 1.85\,\mu\text{m}$, the coherent radiation pulse reaches a peak field amplitude $E_{\text{peak}} = 3.18 \times 10^{11}\text{ V/m}$. Comparing this amplitude against the Sauter-Schwinger vacuum critical field $E_{\text{crit}} = 1.323 \times 10^{18}\text{ V/m}$ yields a field ratio of $E_{\text{peak}} / E_{\text{crit}} = 2.40 \times 10^{-7}$, confirming that operations remain six orders of magnitude below the threshold for non-linear quantum electrodynamic vacuum breakdown and spontaneous electron-positron pair creation.\end{sloppypar}
